\section*{Exercises on \S\arabic{exercisegroup}}
\phantomsection
\addcontentsline{toc}{section}{Exercises on \protect\S\arabic{exercisegroup}}

\begin{enumerate}
\def\labelenumi{\arabic{enumi})}
\tightlist
\item
  Establish the formulae
\end{enumerate}

\[
\tan z = \sum_ {n = 1} ^ {\infty} (- 1) ^ {n - 1} 2 ^ {2 n} (2 ^ {2 n} - 1) b _ {2 n} \frac {z ^ {2 n - 1}}{(2 n) !}
\]

\[
\frac {1}{\sin z} = \frac {1}{z} + \sum_ {n = 1} ^ {\infty} (- 1) ^ {n - 1} 2 (2 ^ {2 n - 1} - 1) b _ {2 n} \frac {z ^ {2 n - 1}}{(2 n) !}
\]

where the series on the right-hand sides converge absolutely, the first for \(|z| < \frac{\pi}{2}\) and the second for \(|z| < \pi\) (express \(\tan z\) and \(1/\sin 2z\) as linear combinations of \(\cot z\) and \(\cot 2z\) ). Deduce that the numbers \(\frac{2^{2n}(2^{2n}-1)}{2n} b_{2n}\) are integers. (Use the following lemma: if, in two absolutely convergent series \(\sum_{n=0}^{\infty} \alpha_n \frac{z^n}{n!}\) , \(\sum_{n=0}^{\infty} \beta_n \frac{z^n}{n!}\) the coefficients \(\alpha_n\) and \(\beta_n\) are integers, then, in their product written in the form \(\sum_{n=0}^{\infty} \gamma_n \frac{z^n}{n!}\) , the \(\gamma_n\) are integers.)

\begin{enumerate}
\def\labelenumi{\arabic{enumi})}
\setcounter{enumi}{1}
\tightlist
\item
  Establish the formula
\end{enumerate}

\[
(n - 1) \mathrm{B} _ {n} (\mathrm{X}) = n (\mathrm{X} - 1) \mathrm{B} _ {n - 1} (\mathrm{X}) - \sum_ {k = 0} ^ {n} \binom {n} {k} b _ {k}   \mathrm{B} _ {n - k} (\mathrm{X})
\]

(differentiate the series \(Se^{SX}/(e^{S}-1)\) with respect to S). Deduce the formula

\[
(2 n + 1) b _ {2 n} = - \sum_ {k = 1} ^ {n - 1} \binom {2 n} {2 k} b _ {2 k}   b _ {2 n - 2 k}
\]

for the Bernoulli numbers.

\begin{enumerate}
\def\labelenumi{\arabic{enumi})}
\setcounter{enumi}{2}
\tightlist
\item
  Show that, for every integer p \textgreater{} 1,
\end{enumerate}

\[
\mathrm{B} _ {n} \left(\frac {x}{p}\right) + \mathrm{B} _ {n} \left(\frac {x + 1}{p}\right) + \dots + \mathrm{B} _ {n} \left(\frac {x + p - 1}{p}\right) = \frac {1}{p ^ {n - 1}} \mathrm{B} _ {n} (\mathrm{X}).
\]

\begin{enumerate}
\def\labelenumi{\arabic{enumi})}
\setcounter{enumi}{3}
\tightlist
\item
  \begin{enumerate}
  \def\labelenumii{\alph{enumii})}
  \tightlist
  \item
    Prove the relation
  \end{enumerate}
\end{enumerate}

\[
\mathrm{B} _ {n} (1 - \mathrm{X}) = (- 1) ^ {n} \mathrm{B} _ {n} (\mathrm{X})
\]

(use the fact that \(b_{2n - 1} = 0\) for \(n > 1\) , and the relation

\[
\mathrm{B} _ {n} (1 - \mathrm{X}) - \mathrm{B} _ {n} (- \mathrm{X}) = (- 1) ^ {n} n \mathrm{X} ^ {n - 1}.
\]

\begin{enumerate}
\def\labelenumi{\alph{enumi})}
\setcounter{enumi}{1}
\tightlist
\item
  Show that
\end{enumerate}

\[
\mathrm{B} _ {n} \left(\frac {1}{2}\right) = b _ {n} \left(\frac {1}{2 ^ {n - 1}} - 1\right)
\]

(use exerc. 3).

\begin{enumerate}
\def\labelenumi{\alph{enumi})}
\setcounter{enumi}{2}
\tightlist
\item
  Show that, for n even, \(\mathrm{B}_{n}(\mathrm{X})\) has two roots in the interval {[}0, 1{]} of R, and that for n odd \textgreater{} 1, \(\mathrm{B}_{n}(\mathrm{X})\) has a simple root at the points 0, \(\frac{1}{2}\) and 1 and does not vanish at any other point of {[}0, 1{]} (use b) and the relation \(B_{n}^{\prime}=nB_{n-1}\) ).
\end{enumerate}

\(d\) ) Deduce from \(c\) ) that, for \(n\) even, the maximum of \(|\mathrm{B}_n(x)|\) on the interval {[}0, 1{]} is \(|b_n|\) , and that for \(n\) odd, if \(a_{n}\) is the maximum of \(|\mathrm{B}_n(x)|\) on {[}0, 1{]}, then

\[
\frac {4}{n + 1} \left| b _ {n + 1} \right| \left(1 - \frac {1}{2 ^ {n}}\right) \leqslant a _ {n} \leqslant \frac {1}{2} n \left| b _ {n - 1} \right|
\]

(use the mean value theorem).

\begin{enumerate}
\def\labelenumi{\arabic{enumi})}
\setcounter{enumi}{4}
\tightlist
\item
  If one puts \(S_{n}(x) = \frac{1}{n + 1} (\mathrm{B}_{n + 1}(x) - \mathrm{B}_{n + 1}(0))\) then, for every integer \(a > 0\) one has
\end{enumerate}

\[
\mathrm{S} _ {n} (a) = 1 ^ {n} + 2 ^ {n} + \dots + (a - 1) ^ {n}.
\]

\begin{enumerate}
\def\labelenumi{\alph{enumi})}
\item
  Show that for every integer \(n \geqslant 0\) and every integer \(a > 0\) one has \(2\mathrm{S}_{2n + 1}(a) \equiv 0 \pmod{a}\) (consider the sum \(k^{2n + 1} + (a - k)^{2n + 1}\) ).
\item
  If \(r\) and \(s\) are any two integers show that
\end{enumerate}

\[
\mathrm{S} _ {n} (r s) \equiv s \mathrm{S} _ {n} (r) + n r \mathrm{S} _ {n - 1} (r) \mathrm{S} _ {1} (s) \pmod {r ^ {2}}.
\]

\begin{enumerate}
\def\labelenumi{\alph{enumi})}
\setcounter{enumi}{2}
\tightlist
\item
  Let p be a prime number. Show that if n is divisible by p-1 one has \(S_{n}(p) \equiv -1 \pmod{p}\) , and if n is not divisible by p-1 then \(S_{n}(p) \equiv 0 \pmod{p}\) (if p does not divide the integer g remark that \(S_{n}(p) \equiv g^{n}S_{n}(p) \pmod{p}\) ).
\end{enumerate}

\begin{enumerate}
\def\labelenumi{\arabic{enumi})}
\setcounter{enumi}{5}
\tightlist
\item
  \begin{enumerate}
  \def\labelenumii{\alph{enumii})}
  \tightlist
  \item
    The rational numbers \(b_{n}\) having been defined by the formula (20) of VI, p.~275, one denotes by \(d_{n}\) the denominator \textgreater{} 0 of \(b_{n}\) written as an irreducible fraction. Show that no prime factor of \(d_{n}\) can be \(> n + 1\) (use the induction formula (23) of VI, p.~276).
  \end{enumerate}
\end{enumerate}

\begin{enumerate}
\def\labelenumi{\alph{enumi})}
\setcounter{enumi}{1}
\tightlist
\item
  Show that for every integer p \textgreater{} 0 and every integer n \textgreater{} 0
\end{enumerate}

\[
\mathrm{S} _ {n} (p) = b _ {n} p + \binom {n} {1} \frac {p}{2} b _ {n - 1} p + \dots + \binom {n} {r} \frac {p ^ {r}}{r + 1} b _ {n - r} p + \dots + \frac {p ^ {n + 1}}{n + 1}.
\]

\begin{enumerate}
\def\labelenumi{\alph{enumi})}
\setcounter{enumi}{2}
\item
  Deduce from b) by recursion on n that, for every prime number p the denominator of \(S_{n}(p) - b_{n}p\) written as an irreducible fraction, is not divisible by p (observe that \(p^{r+1}\) cannot divide \(r + 1\) ).
\item
  Deduce from c) that the number
\end{enumerate}

\[
b _ {n} - \sum_ {p} \frac {\mathrm{S} _ {n} (p)}{p}
\]

where \(p\) runs through the set of prime numbers \(p \leqslant n + 1\) and \(n\) is even, is an integer. Conclude that

\[
b _ {2 n} + \sum_ {p} \frac {1}{p}
\]

where \(p\) runs through the set of prime numbers such that \(p - 1\) divides \(2n\) , is an integer (Clausen-von Staudt theorem; use exerc. 5 c).

* 7) We accept that for every integer \(a > 0\) there are infinitely many prime numbers in the set of integers \(1 + ma\) ( \(m\) running through the set of integers \(\geqslant 1\) ; this is a particular case of Dirichlet's theorem on arithmetic progressions).

\begin{enumerate}
\def\labelenumi{\alph{enumi})}
\item
  Let n be an integer \(\geqslant 1\) , and let \(s \geqslant 1\) be an integer such that \(q = 1 + (2n + 1)!s\) is prime; show that if p is a prime number such that p - 1 divides 2nq then p - 1 must divide 2n (in the opposite case one would have p - 1 = qd with d an integer, and p would be divisible by \(d + 1\) ).
\item
  Deduce from a) that for every integer n \textgreater{} 0 there exist infinitely many integers m \textgreater{} n such that \(b_{2m} - b_{2n}\) is an integer.*
\end{enumerate}

\begin{enumerate}
\def\labelenumi{\arabic{enumi})}
\setcounter{enumi}{7}
\item
  Show that, for every prime number \(p > 3\) , \(\mathrm{S}_{2n}(p^k) - p^k b_{2n}\) , written as an irreducible fraction, has a numerator divisible by \(p^{2k}\) (argue as in exerc. 6).
\item
  To say that a rational number \(r\) is a \(p\) -adic integer (Gen.~Top., III, p.~322, exerc. 23) for \(p\) a prime number, means that when \(r\) is expressed as an irreducible fraction its denominator is not divisible by \(p\) ; one writes \(r \equiv 0 \pmod{p}\) to express the fact that \(r/p\) is a \(p\) -adic integer, and \(r \equiv r' \pmod{p}\) is equivalent by definition to \(r - r' \equiv 0 \pmod{p}\) .
\end{enumerate}

\begin{enumerate}
\def\labelenumi{\alph{enumi})}
\tightlist
\item
  Let \(m\) be a rational integer; the function \(F_{m}(z) = \frac{m}{e^{mz} - 1} - \frac{1}{e^{z} - 1}\) is analytic for \(|z|\) sufficiently small, so can be written as an entire series convergent on a neighbourhood of 0
\end{enumerate}

\[
\mathrm{F} _ {m} (z) = \sum_ {n = 1} ^ {\infty} (m ^ {n} - 1) \frac {b _ {n}}{n} \frac {z ^ {n - 1}}{(n - 1) !}.
\]

Show that also on a neighbourhood of 0

\[
\mathrm{F} _ {m} (z) = \sum_ {n = 0} ^ {\infty} c _ {n} (e ^ {z} - 1) ^ {n}
\]

where the \(c_{n}\) are \(p\) -adic integers; deduce that the numbers \(a_{n} = (m^{n} - 1)\frac{b_{n}}{n}\) are \(p\) -adic integers; further, one has \(a_{n + p - 1} \equiv a_{n} (\mathrm{mod.} p)\) .

\[
\frac {b _ {n + p - 1}}{n + p - 1} \equiv \frac {b _ {n}}{n} \pmod {p}
\]

\begin{enumerate}
\def\labelenumi{\alph{enumi})}
\setcounter{enumi}{1}
\tightlist
\item
  Deduce from a) that if \(p - 1\) does not divide \(n\) , then \(b_{n} / n\) is a \(p\) -adic integer and that (Kummer's congruences). (Take for \(m\) an integer whose class mod. \(p\) is a generator of the multiplicative group of the invertible elements of \(\mathbf{Z} / p\mathbf{Z}\) (Alg., VII. 12).)
\end{enumerate}

\begin{enumerate}
\def\labelenumi{\arabic{enumi})}
\setcounter{enumi}{9}
\item
  With the notation of exerc. 9 show that the numbers \(m^n a_n = m^n (m^n - 1)b_n / n\) are integers (write \(F_m(z)\) as the quotient of two polynomials in \(e^z\) ). Deduce that (for \(n\) even \(\geqslant 2\) ) the denominator \(\delta_n\) of \(b_n / n\) in irreducible form is such that, for every integer \(m\) prime to \(n\) one has \(m^n \equiv 1\) (mod. \(\delta_n\) ). Further, if an integer \(d\) is such that \(m^n \equiv 1\) (mod. \(d\) ) for every integer \(m\) prime to \(n\) , then \(d\) divides \(\delta_n\) (use the Clausen-von Staudt theorem and the structure of the multiplicative group of \(\mathbf{Z} / d\mathbf{Z}\) (Alg., VII. 12)).
\item
  \begin{enumerate}
  \def\labelenumii{\alph{enumii})}
  \tightlist
  \item
    Let p be a prime number \(\neq 2\) . Then the following properties are equivalent:
  \end{enumerate}
\end{enumerate}

\(\alpha) b_n / n \not\equiv 0 \pmod{p}\) for every even \(n\) such that \(p - 1\) does not divide \(n\) .

\(\beta) b_n \not\equiv 0 \pmod{p}\) for \(n = 2, 4, 6, \ldots, p - 3\) .

(Use exerc. 9b).)

\begin{enumerate}
\def\labelenumi{\alph{enumi})}
\setcounter{enumi}{1}
\tightlist
\item
  One says that \(p\) is a regular prime number if it is equal to 2 or satisfies the equivalent conditions of \(a\) ); otherwise \(p\) is called irregular. The prime numbers 2, 3, 5, 7, 11 are regular; 691 is irregular.
\end{enumerate}

Let \(I\) be a finite set of irregular primes. Let \(n\) be an integer \(\geqslant 2\) multiple of \(\prod_{q \in I} (q - 1)\)

and such that \(|b_n / n| > 1\) , (cf.~VI, p.~288, formula (15)). Let \(p\) be a prime factor of the numerator of \(|b_n / n|\) in irreducible form. Show that \(p - 1\) does not divide \(n\) and that consequently \(p\) is irregular. Deduce that the set of irregular prime numbers is infinite.

\begin{enumerate}
\def\labelenumi{\arabic{enumi})}
\setcounter{enumi}{11}
\tightlist
\item
  For every polynomial \(\mathbf{Q}\) with real coefficients such that \(\mathrm{Q}(0) = \mathrm{Q}(1)\) one writes \(\tilde{\mathbf{Q}}\) for the continuous function on \(\mathbf{R}\) such that \(\tilde{\mathbf{Q}}(x) = \mathbf{Q}(x)\) for \(0 \leqslant x \leqslant 1\) and \(\tilde{\mathbf{Q}}(x + 1) = \tilde{\mathbf{Q}}(x)\) for every \(x \in \mathbf{R}\) .
\end{enumerate}

\begin{enumerate}
\def\labelenumi{\alph{enumi})}
\item
  For every integer \(m \geqslant 2\) the Bernoulli polynomial \(B_{m}\) is the unique polynomial Q with real coefficients, of degree m, monic, and such that \(\int_{0}^{1} Q(x) dx = 0\) , \(Q(1) = Q(0)\) and such that the function \(\tilde{Q}\) is m - 2 times differentiable on R and has a continuous \((m - 2)^{nd}\) derivative. (Argue by recursion on m.)
\item
  Show that for \(m \geqslant 1\)
\end{enumerate}

\[
\int_ {0} ^ {1} \mathrm{B} _ {m} (x) e ^ {- 2 \pi i n x} d x = \left\{ \begin{array}{c c c} 0 & \text {if} & n = 0 \\ - \frac {m !}{(2 \pi i n) ^ {m}} & \text {if} & n \neq 0 \end{array} \right. \text {on} \mathbf {Z}.
\]

* c) Deduce from b) that, for \(0 \leqslant x \leqslant 1\) and \(m \geqslant 2\) one has

\[
\mathrm{B} _ {m} (x) = - \frac {m !}{(2 \pi i) ^ {m}} \sum_ {n \in \mathbf {Z} - \{0 \}} \frac {e ^ {2 \pi i n x}}{n ^ {m}}
\]

the series being normally convergent.*

\begin{enumerate}
\def\labelenumi{\arabic{enumi})}
\setcounter{enumi}{12}
\tightlist
\item
  Let \(f\) be an integer \(\geqslant 1\) , and \(\chi\) a function defined on \(\mathbf{Z}\) , with complex values, such that \(\chi(x + f) = \chi(x)\) for every \(x \in \mathbf{Z}\) . If one puts
\end{enumerate}

\[
\hat {\chi} (y) = \frac {1}{f} \sum_ {x = 0} ^ {f - 1} \chi (x) e ^ {- 2 \pi i x y / f} \quad \text { for } y \in \mathbf {Z},
\]

one has

\[
\chi (x) = \sum_ {y = 0} ^ {f - 1} \hat {\chi} (y) e ^ {- 2 \pi i x y / f} \quad \text {   for   } x \in \mathbf {Z}.
\]

Deduce from exerc. 12 that, for every integer \(m \geqslant 2\) ,

\[
\sum_ {n \in \mathbf {Z} \backslash \{0 \}} \frac {\chi (n)}{n ^ {m}} = - \frac {(2 \pi i) ^ {m}}{m !} \sum_ {y = 0} ^ {f - 1} \hat {\chi} (y) \mathrm{B} _ {m} \left(\frac {y}{f}\right).
\]

For example, one has

\[
1 - \frac {1}{3 ^ {3}} + \frac {1}{5 ^ {3}} - \dots + \frac {(- 1) ^ {n}}{(2 n + 1) ^ {3}} + \dots = \frac {\pi^ {3}}{3 2}.
\]

\setcounter{exercisegroup}{2}
\refstepcounter{exercisegroup}
